\section{Category similarity}\label{sec:category_similarity}

\begin{remark}\label{rem:category_similarity}
  We have several notions for expressing that two categories \( \cat{C} \) and \( \cat{D} \) are similar:

  \begin{thmenum}
    \thmitem{rem:category_similarity/equality} Obviously, if \( \cat{C} \) and \( \cat{D} \) are equal, they are similar.

    \thmitem{rem:category_similarity/isomorphism} A weaker notion is that of isomorphism, defined in \cref{def:isomorphism_of_categories} as the existence of a fully invertible functor \( F: \cat{C} \to \cat{D} \).

    \thmitem{rem:category_similarity/equivalence} An even weaker but useful notion is that of equivalence, defined in \cref{def:equivalence_of_categories} as a generalization of invertible isomorphisms --- a pair of opposite functors whose compositions are \hyperref[def:natural_isomorphism]{naturally isomorphic} to the corresponding identities.
  \end{thmenum}
\end{remark}

\paragraph{Isomorphism of categories}

\begin{definition}\label{def:isomorphism_of_categories}\mimprovised
  We say that the \hyperref[def:functor]{functor} \( F: \cat{C} \to \cat{D} \) between arbitrary (possibly \hyperref[def:small_category]{large}) categories is an \term{isomorphism} if and only if it has a two-sided inverse; that is, a functor \( G: \cat{D} \to \cat{C} \) such that
  \begin{align*}
    G \bincirc F &= \id_{\cat{C}}, \\
    F \bincirc G &= \id_{\cat{D}}.
  \end{align*}

  If an isomorphism between \( \cat{C} \) and \( \cat{D} \) exists, we say that the categories are \term{isomorphic}.
\end{definition}
\begin{comments}
  \item This is a generalization of isomorphisms in the category \hyperref[def:category_of_small_categories]{\( \cat{Cat} \)} of small categories.
\end{comments}

\begin{proposition}\label{thm:isomorphism_of_categories_iff_bijective_on_morphisms}
  The \hyperref[def:functor]{functor} \( F: \cat{C} \to \cat{D} \) is \hyperref[def:functor_invertibility/bijective]{bijective on morphisms} if and only if it is an \hyperref[def:isomorphism_of_categories]{isomorphism of categories}.
\end{proposition}
\begin{proof}
  This proposition generalizes \cref{thm:small_functor_invertibility/two_sided}.
\end{proof}

\begin{example}\label{ex:def:isomorphism_of_categories}
  We list examples of \hyperref[def:isomorphism_of_categories]{isomorphisms of categories}:
  \begin{thmenum}
    \thmitem{ex:def:isomorphism_of_categories/discrete_cat} Denote by \( \cat{DCat} \) the category of all small \hyperref[def:discrete_category]{discrete categories} and functors between them.

    Consider the \hyperref[con:forgetful_functor]{forgetful functor} \( U: \cat{DCat} \to \cat{Set} \) that sends a category to its set of objects. The discrete category functor \( D: \cat{Set} \to \cat{DCat} \) is its obvious two-sided inverse.

    Thus, \( \cat{DCat} \) and \( \cat{Set} \) are actually isomorphic categories.

    \thmitem{ex:def:isomorphism_of_categories/preorder} As shown in \fullref{thm:order_category_isomorphism}, the category of small \hyperref[def:preordered_set]{preordered sets} is isomorphic to that of small \hyperref[def:thin_category]{thin categories}.

    \thmitem{ex:def:isomorphism_of_categories/integers} The following isomorphisms are are closely related:
    \begin{itemize}
      \item \Cref{thm:commutative_monoid_is_semimodule}: \hyperref[def:monoid]{commutative monoids} and \hyperref[def:semimodule]{\( \BbbN \)-semimodules}.

      \item \Cref{thm:semiring_is_natural_number_algebra}: \hyperref[def:semiring]{semirings} and \hyperref[def:algebra_over_semiring]{associative \( \BbbN \)-algebras}.

      \item \Cref{thm:abelian_group_is_module}: \hyperref[def:abelian_group]{abelian groups} and \hyperref[def:module]{\( \BbbZ \)-modules}.

      \item \Cref{thm:ring_is_integer_algebra}: \hyperref[def:ring]{rings} and \hyperref[def:algebra_over_semiring]{associative \( \BbbZ \)-algebras}.
    \end{itemize}
  \end{thmenum}
\end{example}

\paragraph{Thin categories}

\begin{definition}\label{def:thin_category}\mcite[example 3.26(3)]{AdámekHerrlichStrecker2004JoyOfCats}
  We call a \hyperref[def:category]{category} \term{thin} if parallel morphisms are equal.
\end{definition}
\begin{comments}
  \item The \( \hom \)-sets of a thin category are \hyperref[def:subsingleton_set]{subsingleton sets}.

  \item Thin categories can be conflated with preordered sets due to \fullref{thm:order_category_isomorphism}. In fact, thin categories are called \enquote{preorders} in \cite[11]{MacLane1998CategoryTheory}.
\end{comments}

\begin{definition}\label{def:skeletal_category}\mcite[93]{MacLane1998CategoryTheory}
  We call a \hyperref[def:category]{category} is \term{skeletal} if \hyperref[def:morphism_invertibility/isomorphism]{isomorphic objects} are equal.
\end{definition}
\begin{comments}
  \item Skeletal categories are used to define skeletons in \cref{def:category_skeleton}.
\end{comments}

\begin{theorem}[Ordered sets as categories]\label{thm:order_category_isomorphism}
  The category of \hyperref[def:small_set]{small} \hyperref[def:preordered_set]{preordered sets} and their homomorphisms is \hyperref[def:isomorphism_of_categories]{isomorphic} to the category of \hyperref[def:small_category]{small} \hyperref[def:thin_category]{thin categories} and their functors.

  \hyperref[def:partially_ordered_set]{Partially ordered sets} correspond to \hyperref[def:skeletal_category]{skeletal} thin categories.
\end{theorem}
\begin{proof}
  \SubProof{Proof that preordered sets induce thin categories} Let \( (P, \leq) \) be a small preordered set. Let \( \cat{P} \) the following category \( \cat{P} \)\fnote{Modulo technicalities, this is the \hyperref[def:free_category]{free category} obtained by regarding \( (P, \leq) \) as a \hyperref[def:directed_multigraph]{directed multigraph}, as per \cref{def:directed_graph_induced_by_relation}.}:
  \begin{itemize}
    \item The \hyperref[def:category/objects]{set of objects} \( \obj(\cat{P}) \) is simply \( P \).

    \item The \hyperref[def:category/morphisms]{set of morphisms} \( \cat{P}(x, y) \) consists of the tuple \( (x, y) \) if \( x \leq y \) and is empty otherwise.

    \item The \hyperref[def:category/composition]{composition of the morphisms} \( (x, y) \) and \( (y, z) \) is simply \( (x, y) \). This is well-defined because of the \hyperref[def:binary_relation/transitive]{transitivity} of \( \leq \).

    \item The \hyperref[def:category/identity]{identity morphism} on an object \( x \) is \( (x, x) \). This is well-defined because of the \hyperref[def:binary_relation/reflexive]{reflexivity} of \( \leq \).
  \end{itemize}

  This category \( \cat{P} \) is thin because ordered tuples with the same elements are equal.

  Now let \( f: P \to Q \) be a nonstrict order-preserving map. It induces the functor
  \begin{equation*}
    \begin{aligned}
      &F: \cat{P} \to \cat{Q} \\
      &F(x) \coloneqq f(x) \\
      &F(x, y) \coloneqq (f(x), f(y)).
    \end{aligned}
  \end{equation*}

  Both functor properties \ref{def:functor/CF1} and \ref{def:functor/CF2} are immediate.

  \SubProof{Proof that thin categories induce preordered sets} Let \( \cat{P} \) be a small category. Define the binary relation
  \begin{equation*}
    X \leq Y \T{if and only if} \cat{P}(X, Y) \neq \varnothing.
  \end{equation*}

  This is a binary relation over the set \( P \coloneqq \obj(\cat{P}) \). It is reflexive because of the existence of identity morphisms in \( \cat{P} \) and transitive because of the requirement that the composition of compatible morphisms exists. Therefore, \( (P, \leq) \) is a preordered set.

  Given a functor \( F: \cat{P} \to \cat{Q} \), the restriction \( F\restr_{\obj(\cat{P})} \) is a nonstrict order-preserving map from \( (P, \leq_P) \) to \( (Q, \leq_Q) \).

  Indeed, if \( X \leq Y \) for objects \( X, Y \) in \( \cat{P} \), then \( \cat{P}(X, Y) \neq \varnothing \). Hence, \( \cat{P}(F(X), F(Y)) \neq \varnothing \) and \( F(X) \leq F(Y) \).

  \SubProof{Proof of category isomorphism} We have implicitly defined mutually inverse functors between the category of small preordered sets and that of small thin categories.

  \SubProof{Proof for partial orders} Fix a preordered set \( (P, \leq) \) and its thin category \( \cat{P} \).

  Suppose first that antisymmetry holds in \( (P, \leq) \), and that \( x \) and \( y \) are isomorphic in \( \cat{P} \). Then \( (x, y) \) and \( (y, x) \) are inverse morphisms in \( \cat{P} \). Thus, \( x \leq y \) and \( y \leq x \) both hold in \( (P, \leq) \). By antisymmetry, we have \( x = y \).

  Conversely, suppose that \( \cat{P} \) is skeletal. If \( x \leq y \) and \( y \leq x \) in \( \cat{P} \), then \( (x, y) \) and \( (y, x) \) are both morphisms in \( \cat{P} \), and are clearly mutually inverse, so \( x = y \) since \( \cat{P} \) is skeletal.

  Therefore, antisymmetry holds in \( (P, \leq) \) if and only if \( \cat{P} \) is skeletal.
\end{proof}

\begin{proposition}\label{thm:order_category_isomorphism_properties}
  Let \( (P, \leq) \) be a \hyperref[def:partially_ordered_set]{partially ordered set} and let \( \cat{P} \) be its corresponding \hyperref[def:thin_category]{thin category} induced by \fullref{thm:order_category_isomorphism}.

  \begin{thmenum}
    \thmitem{thm:order_category_isomorphism_properties/opposite} The \hyperref[def:opposite_category]{opposite category} \( \cat{P}^\oppos \) corresponds to the \hyperref[def:partially_ordered_set/opposite]{opposite ordered set} \( (P, \geq) \).

    \thmitem{thm:order_category_isomorphism_properties/universal} The \hyperref[def:universal_objects/terminal]{terminal} and \hyperref[def:universal_objects/initial]{initial objects} in \( \cat{P} \) correspond to (possibly nonunique) \hyperref[def:extremal_points/top_and_bottom]{top and bottom}\fnote{In this context, a top (resp. bottom) element is a global upper (resp. loewr) bound.} elements in \( (P, \leq) \), respectively.
  \end{thmenum}
\end{proposition}
\begin{proof}
  \SubProofOf{thm:order_category_isomorphism_properties/opposite} Trivial.

  \SubProofOf{thm:order_category_isomorphism_properties/universal} Trivial.
\end{proof}

\paragraph{Equivalences of categories}

\begin{definition}\label{def:equivalence_of_categories}\mcite[def. 1.3.15]{Awodey2010CategoryTheory}
  An \term{equivalence} of the \hyperref[def:category]{categories} \( \cat{C} \) and \( \cat{D} \) is a quadruple
  \begin{equation*}
    \begin{aligned}
                F &: \cat{C} \to \cat{D}, \\
                G &: \cat{D} \to \cat{C}, \\
             \eta &: \id_{\cat{C}} \Rightarrow G \bincirc F, \\
      \varepsilon &: F \bincirc G \Rightarrow \id_{\cat{D}},
    \end{aligned}
  \end{equation*}
  where \( \eta \) and \( \varepsilon \) are \hyperref[thm:natural_isomorphism]{natural isomorphisms}.

  If \( (F, G, \eta, \varepsilon) \) is an equivalence, so is \( (G, F, \varepsilon^{-1}, \eta^{-1}) \), so the roles of \( \cat{C} \) and \( \cat{D} \) are symmetric if we are only interested in whether an equivalence between them exists. In this case wee say that \( \cat{C} \) and \( \cat{D} \) are equivalent categories.
\end{definition}
\begin{comments}
  \item As can be expected, equivalence of categories is an equivalence relation on the objects of \( \cat{Cat} \) --- see \cref{thm:category_equivalence_is_equivalence_relation}.

  \item This definition is adjusted so that \( \eta \) and \( \varepsilon \) have the same signature as their eponymous natural transformations in \hyperref[def:unit_counit_adjunction]{unit-counit adjunctions}.

  Note that an equivalence is not necessarily an adjunction, they simply have a common setup. An \hyperref[def:adjoint_equivalence]{adjoint equivalence} combines the notions; \cref{thm:adjoint_equivalence} shows how every equivalence can be converted to an adjoint equivalence.
\end{comments}

\begin{example}\label{ex:def:equivalence_of_categories}
  We list examples of \hyperref[def:equivalence_of_categories]{equivalences of categories}:
  \begin{thmenum}
    \thmitem{ex:def:equivalence_of_categories/groupoid} \Cref{thm:connected_delooping} shows how a \hyperref[def:connected_category]{connected} \hyperref[def:groupoid]{groupoid} determines a \hyperref[def:group]{group} up to an equivalence of categories.

    \thmitem{ex:def:equivalence_of_categories/skeleton} \Cref{thm:category_skeleton_existence} shows how every category is equivalence to each of its \hyperref[def:category_skeleton]{skeletons}.
  \end{thmenum}
\end{example}

\begin{proposition}\label{thm:category_equivalence_is_equivalence_relation}
  \hyperref[def:equivalence_of_categories]{Category equivalence} is an \hyperref[def:equivalence_relation]{equivalence relation} on the set of all \hyperref[def:small_category]{small categories} (the set of objects of \hyperref[def:category_of_small_categories]{\( \cat{Cat} \)}).
\end{proposition}
\begin{proof}
  \SubProofOf[def:binary_relation/reflexive]{reflexivity} Clearly \( (\id_{\cat{C}}, \id_{\cat{C}}, \id_{\id_{\cat{C}}}, \id_{\id_{\cat{C}}}]) \) is a self-equivalence for the category \( \cat{C} \).

  \SubProofOf[def:binary_relation/symmetric]{symmetry} If \( (F, G, \eta, \varepsilon) \) is an equivalence, so is \( (G, F, \varepsilon^{-1}, \eta^{-1}) \).

  \SubProofOf[def:binary_relation/transitive]{transitivity} By \cref{thm:def:morphism_invertibility/composition}, the composition of invertible natural isomorphisms is again a natural isomorphism, hence equivalence in \( \obj(\cat{Cat}) \) is a transitive relation.
\end{proof}

\begin{proposition}\label{thm:category_isomorphism_is_equivalence}
  \hyperref[def:isomorphism_of_categories]{Isomorphic categories} are \hyperref[def:equivalence_of_categories]{equivalent}.

  More precisely, if \( F: \cat{C} \to \cat{D} \) is an isomorphism with an inverse \( G: \cat{D} \to \cat{C} \), by definition
  \begin{align*}
    G \bincirc F &= \id_{\cat{C}}, \\
    F \bincirc G &= \id_{\cat{D}}.
  \end{align*}

  Then \( (F, G, \id_{\id_{\cat{C}}}, \id_{\id_{\cat{D}}}) \) is an equivalence.
\end{proposition}
\begin{proof}
  Trivial.
\end{proof}

\begin{proposition}\label{thm:discrete_category_equivalence}
  The \hyperref[def:discrete_category]{discrete categories} \( \cat{C} \) and \( \cat{D} \) are \hyperref[def:equivalence_of_categories]{equivalent} if and only if the underlying sets \( \obj(\cat{C}) \) and \( \obj(\cat{D}) \) are \hyperref[def:equinumerosity]{equinumerous}.
\end{proposition}
\begin{proof}
  \SufficiencySubProof Suppose that \( (F, G, \eta, \varepsilon) \) is a category equivalence.

  Since \( \cat{D} \) is discrete, the only morphisms in \( \cat{D} \) are the identities; thus, the natural transformation \( \varepsilon \) consists of identity morphisms. This in turn implies that \( F \bincirc G = \id_{\cat{D}} \).

  We similarly conclude that \( G \bincirc F = \id_{\cat{C}} \), and hence that \( F \) is an \hyperref[def:isomorphism_of_categories]{isomorphism of categories}. \Cref{thm:isomorphism_of_categories_iff_bijective_on_morphisms} implies that \( F \) is bijective on morphisms, and \cref{thm:def:functor_invertibility/bijective} implies that \( F \) is bijective on objects.

  Therefore, the sets \( \obj(\cat{C}) \) and \( \obj(\cat{D}) \) are equinumerous.

  \NecessitySubProof Suppose that \( f: \obj(\cat{C}) \to \obj(\cat{D}) \) is a bijective function. Let \( F: \cat{C} \to \cat{D} \) be the functor that extends \( f \) to identity morphisms. Similarly, let \( G: \cat{D} \to \cat{C} \) be an extension of \( f^{-1} \). Then \( G \bincirc F = \id_{\cat{C}} \) and \( F \bincirc G = \id_{\cat{D}} \), i.e. \( F \) is an isomorphism.

  By \cref{thm:category_isomorphism_is_equivalence}, \( (F, G, \id_{\id_{\cat{C}}}, \id_{\id_{\cat{D}}}) \) is an equivalence.
\end{proof}

\begin{proposition}\label{thm:opposite_of_category_equivalence}
  The \hyperref[def:opposite_category]{opposites} of \hyperref[def:equivalence_of_categories]{equivalent categories} are equivalent.

  More precisely, fix an equivalence \( (F, G, \eta, \varepsilon) \) between \( \cat{C} \) and \( \cat{D} \). Consider the quadruple
  \begin{equation*}
    \begin{aligned}
                G^\oppos &: \cat{D}^\oppos \to \cat{C}^\oppos, \\
                F^\oppos &: \cat{C}^\oppos \to \cat{D}^\oppos, \\
      \varepsilon^\oppos &: \id_{\cat{D}} \Rightarrow (F \bincirc G)^\oppos, \\
             \eta^\oppos &: \underbrace{(G \bincirc F)^\oppos}_{G^\oppos \bincirc F^\oppos} \Rightarrow \id_{\cat{C}^\oppos},
    \end{aligned}
  \end{equation*}
  where \( G^\oppos \) and \( F^\oppos \) are \hyperref[def:opposite_functor]{opposite functors} and \( \varepsilon^\oppos \) and \( \eta^\oppos \) are \hyperref[def:opposite_natural_transformation]{opposite natural transformations}. We claim that this is an equivalence between \( \cat{D}^\oppos \) and \( \cat{C}^\oppos \).
\end{proposition}
\begin{comments}
  \item \Cref{thm:composition_of_opposite_functors} implies that \( (F \bincirc G)^\oppos = F^\oppos \bincirc G^\oppos \) and \( (G \bincirc F)^\oppos = G^\oppos \bincirc F^\oppos \).

  \item This is part of the categorical duality principles listed in \cref{ex:thm:category_duality}.
\end{comments}
\begin{proof}
  Trivial.
\end{proof}

\begin{theorem}[Characterization of equivalences of categories]\label{thm:categorical_equivalence_characterization}
  A \hyperref[def:functor]{functor} \( F: \cat{C} \to \cat{D} \) is \hyperref[def:functor_invertibility/fully_faithful]{fully faithful} and \hyperref[def:functor_invertibility/essentially_surjective]{essentially surjective} if and only if there exists a functor \( G: \cat{D} \to \cat{C} \) and \hyperref[def:natural_isomorphism]{natural isomorphisms}
  \begin{align*}
    \eta        &: \id_{\cat{C}} \Rightarrow G \bincirc F, \\
    \varepsilon &: F \bincirc G \Rightarrow \id_{\cat{D}},
  \end{align*}
\end{theorem}
\begin{comments}
  \item Such a quadruple \( (F, G, \eta, \varepsilon) \) is by definition an \hyperref[def:equivalence_of_categories]{equivalence of categories}.

  \item In \hyperref[def:zfc]{\logic{ZF}}, this theorem is equivalent to the \hyperref[def:zfc/choice]{axiom of choice}. We list equivalent results in \cref{thm:axiom_of_choice_equivalences}.
\end{comments}
\begin{proof}
  \SubProof{Proof of sufficiency assuming the axiom of choice} Suppose that the axiom of choice holds and let \( F \) be a fully faithful functor that is surjective on objects.

  \SubProof*{Construction of \( G \)} From essential surjectivity of \( F \), it follows that for every object \( X \) in \( \cat{D} \), the \enquote{essential preimage} of \( X \) under \( F \), i.e. the set \( \mscrA_X \) of objects \( A \) in \( \cat{C} \) such that \( F(A) \cong X \), is nonempty.

  We use the axiom of choice on the family \( \set{ \mscrA_X }_{X \in \obj(\cat{D})} \) to select a single preimage for every \( X \). This gives us a function \( G_{\obj}: \obj(\cat{D}) \to \obj(\cat{C}) \).

  There may be multiple isomorphisms from \( F(G_{\obj}(X)) \) to \( X \). We use the axiom of choice to pick a particular an isomorphism \( \varepsilon_X: F(G_{\obj}(X)) \to X \) for every \( X \).

  In order to \( G \) to become a functor, we must extend it to morphisms. Fix a morphism \( g: X \to Y \) in \( \cat{D} \) and let \( g' \) be the composite
  \begin{equation*}
    \begin{aligned}
      \includegraphics[page=1]{output/thm__categorical_equivalence_characterization}
    \end{aligned}
  \end{equation*}

  This morphism \( g' \) belongs to the \( \hom \)-set
  \begin{equation*}
    \cat{D}(F(G_{\obj}(X)), F(G_{\obj}(Y))).
  \end{equation*}

  Since \( F \) is fully faithful, there exists a unique morphism \( f: G_{\obj}(X) \to G_{\obj}(Y) \) such that \( F(f) = g' \). We put \( G_{\mor}(g) \coloneqq f \).

  Thus defined, \( G \) satisfies
  \begin{equation*}
    \begin{aligned}
      \includegraphics[page=2]{output/thm__categorical_equivalence_characterization}
    \end{aligned}
  \end{equation*}

  This diagram demonstrates the naturality of \( \varepsilon \).

  \SubProof*{Proof of functoriality of \( G \)} To prove \ref{def:functor/CF1} for \( G \), we must first note that, for \( g = \id_X \), naturality of \( \varepsilon \) implies
  \begin{equation*}
    [F \bincirc G](\id_X) = \varepsilon_X^{-1} \bincirc \id_X \bincirc \varepsilon_X.
  \end{equation*}

  The right side equals
  \begin{equation*}
    \id_{[F \bincirc G]}(X)
    \reloset {\ref{def:functor/CF1}} =
    F(\id_{G(X)}).
  \end{equation*}

  Both \( G(\id_X) \) and \( \id_{G(X)} \) belong to \( \cat{D}([F \bincirc G](X), [F \bincirc G](X)) \). Since \( F \) is faithful, we conclude that \( G(\id_X) = \id_{G(X)} \), as desired.

  To prove \ref{def:functor/CF2}, we use the naturality of \( \varepsilon \) for \( g: X \to Y \) and \( h: Y \to Z \):
  \begin{equation*}
    \begin{aligned}
      \includegraphics[page=3]{output/thm__categorical_equivalence_characterization}
    \end{aligned}
  \end{equation*}

  Thus,
  \begin{equation*}
    [F \bincirc G](h) \bincirc [F \bincirc G](g) = \varepsilon_Z^{-1} \bincirc h \bincirc g \bincirc \varepsilon_X = [F \bincirc G](h \bincirc g).
  \end{equation*}

  By the same uniqueness argument as above, we conclude that
  \begin{equation*}
    G(h \bincirc g) = G(h) \bincirc G(g).
  \end{equation*}

  \SubProof*{Construction of \( \eta \)} For every object \( A \) in \( \cat{C} \), we have an isomorphism
  \begin{equation*}
    \varepsilon_{F(A)}^{-1}: F(A) \to [F \bincirc G \bincirc F](A).
  \end{equation*}

  Denote by \( \eta_A \) the unique morphism from \( \cat{C}\parens[\big]{ A, [F \bincirc G](A) } \) such that \( F(\eta_A) = \varepsilon_{F(A)}^{-1} \). This morphism is well-defined because \( F \) is fully faithful by assumption.

  By \cref{thm:def:functor_invertibility/fully_faithful_reflects_invertible}, since \( \varepsilon_{F(A)}^{-1} \) is an isomorphism, so is \( \eta_A \).

  It remains to show that the family \( \eta \) of all such morphisms is a natural transformation. By the naturality of \( \varepsilon^{-1} \), the following diagram commutes:
  \begin{equation*}
    \begin{aligned}
      \includegraphics[page=4]{output/thm__categorical_equivalence_characterization}
    \end{aligned}
  \end{equation*}

  Hence, by \cref{thm:commutative_diagrams_preserved_and_reflected}, the following diagram also commutes:
  \begin{equation*}
    \begin{aligned}
      \includegraphics[page=5]{output/thm__categorical_equivalence_characterization}
    \end{aligned}
  \end{equation*}

  \SubProof{Proof of necessity assuming the axiom of choice} Suppose that \( (F, G, \eta, \varepsilon) \) is an equivalence.

  \SubProofOf*[def:functor_invertibility/essentially_surjective]{essential surjectivity} For any object \( X \) in \( \cat{D} \), \( A \coloneqq G(X) \) is an object in \( \cat{C} \).

  By definition of category equivalence, the morphism
  \begin{equation*}
    \varepsilon_X: \underbrace{[F \bincirc G](X)}_{F(A)} \to X
  \end{equation*}
  is an isomorphism.

  Therefore, for every object \( X \) in \( \cat{D} \), there exists some object \( A \) in \( \cat{C} \) such that \( F(A) \cong X \). Thus, \( F \) is essentially surjective.

  \SubProofOf*[def:functor_invertibility/faithful]{faithfulness} Fix some objects \( A \) and \( B \) in \( \cat{C} \). Let \( f_1: A \to B \) and \( f_2: A \to B \) be morphisms such that \( F(f_1) = F(f_2) \).

  From the naturality of \( \eta \) it follows that the following diagram commutes:
  \begin{equation}\label{eq:thm:equivalence_induces_fully_faithful_and_essentially_surjective_functor/faithfullness}
    \begin{aligned}
      \includegraphics[page=6]{output/thm__categorical_equivalence_characterization}
    \end{aligned}
  \end{equation}

  Therefore, \( \eta_B \bincirc f_1 = \eta_B \bincirc f_2 \). By \cref{thm:def:morphism_invertibility/split}, \( \eta_B \) is left cancellable, thus \( f_1 = f_2 \).

  \SubProofOf*[def:functor_invertibility/full]{fullness} Fix some objects \( A \) and \( B \) in \( \cat{C} \). Let \( g: F(A) \to F(B) \) be an arbitrary morphism.

  We can define a morphism \( f: A \to B \) as the composition
  \begin{equation*}
    \begin{aligned}
      \includegraphics[page=7]{output/thm__categorical_equivalence_characterization}
    \end{aligned}
  \end{equation*}

  By the naturality of \( \eta \) and \( \eta^{-1} \), the following diagram commutes:
  \begin{equation*}
    \begin{aligned}
      \includegraphics[page=8]{output/thm__categorical_equivalence_characterization}
    \end{aligned}
  \end{equation*}

  Therefore,
  \begin{equation*}
    G(g) = [G \bincirc F](f).
  \end{equation*}

  We have already shown that \( F \) is faithful. Using the same argument for the equivalence \( (G, F, \varepsilon^{-1}, \eta^{-1}) \), we can conclude that \( G \) is also faithful. Then, since \( g \) and \( F(f) \) are parallel, we can cancel \( G \) and conclude that \( g = F(f) \).

  Therefore, \( F \) is full.

  \ImplicationSubProof[thm:categorical_equivalence_characterization]{this characterization}[def:zfc/choice]{the axiom of choice} Suppose that fully faithful and essentially surjective functors induce equivalences.

  Let \( \mscrA \) be a family of nonempty sets. Let \( \cat{A} \) be the \hyperref[def:discrete_category]{discrete category} induced by \( \mscrA \).

  Define the category \( \cat{C} \) as follows:
  \begin{itemize}
    \item The \hyperref[def:category/objects]{set of objects} \( \obj(\cat{C}) \) is the \hyperref[def:disjoint_union]{disjoint union} \( \coprod_{A \in \mscrA} A \).

    \item The \hyperref[def:category/morphisms]{set of morphisms} \( \cat{C}(\iota_A(x), \iota_B(y)) \) has a single morphism if \( A = B \) and no morphisms otherwise. This single morphism can be encoded as the triple \( (A, x, y) \).

    \item The \hyperref[def:category/composition]{composition of the morphisms} \( (A, x, y) \) and \( (A, y, z) \) is the morphism \( (A, x, z) \).

    \item The \hyperref[def:category/identity]{identity morphism} of \( (A, x) \) is \( (A, x, x) \).
  \end{itemize}

  Define the functor
  \begin{equation*}
    \begin{aligned}
      &F: \cat{C} \to \cat{A}, \\
      &F(\iota_A(x)) \coloneqq A, \\
      &F(A, x, y) \coloneqq \id_A.
    \end{aligned}
  \end{equation*}

  For each set \( A \) in \( \mscrA \), \( F \) maps the inclusion \( \iota_A(x) \) of a point \( x \) in \( A \) to \( A \). We have taken the disjoint union of \( \mscrA \) since otherwise there may not be a canonical choice of set \( A \) for \( F \) to send \( x \) to. Thus, the functor is surjective on objects (not essentially surjective but actually surjective).

  Note that \( \cat{A}(F(\iota_A(x)), F(\iota_B(y))) \) has a single morphism if \( A = B \) and is empty otherwise. From this it follows that \( F \) is fully faithful.

  Therefore, \( F \) induces an equivalence \( (F, G, \eta, \varepsilon) \). The functor \( G \) chooses an object \( \iota_A(x) \) of \( \cat{C} \) for each object \( A \) of \( \cat{A} \). This induces a \hyperref[def:choice_function]{choice function} on \( \mscrA \).

  We have shown that the axiom of choice holds.
\end{proof}

\paragraph{Category skeletons}

\begin{definition}\label{def:category_skeleton}
  We call a \hyperref[def:subcategory/full]{full} \hyperref[def:subcategory]{subcategory} \( \cat{S} \) of \( \cat{C} \) is a \term[en=skeleton (\cite[93]{MacLane1998CategoryTheory})]{skeleton} if any of the following equivalent conditions hold:

  \begin{thmenum}
    \thmitem{def:category_skeleton/direct}\mcite[93]{MacLane1998CategoryTheory} Every object in \( \cat{C} \) is isomorphic to exactly one object of \( \cat{S} \).

    \thmitem{def:category_skeleton/equivalence}\mcite[def. 16.3.1]{Schubert1972Categories} \( \cat{S} \) is \hyperref[def:skeletal_category]{skeletal} and the inclusion functor \( \Iota: \cat{S} \to \cat{C} \) is \hyperref[def:functor_invertibility/essentially_surjective]{essentially surjective}.
  \end{thmenum}
\end{definition}
\begin{comments}
  \item By \fullref{thm:categorical_equivalence_characterization}, \( \cat{S} \) is \hyperref[def:equivalence_of_categories]{equivalent} to \( \cat{C} \).
\end{comments}
\begin{defproof}
  \ImplicationSubProof{def:category_skeleton/direct}{def:category_skeleton/equivalence} Suppose that that every object in \( \cat{C} \) is isomorphic to exactly one object in \( \cat{S} \).

  Existence in the latter condition implies that \( \Iota \) is essentially surjective, while uniqueness implies that \( \cat{S} \) is skeletal.

  \ImplicationSubProof{def:category_skeleton/equivalence}{def:category_skeleton/direct} Conversely, if \( \Iota: \cat{S} \to \cat{C} \) is essentially surjective, every object \( X \) in \( \cat{C} \) is isomorphic to some \( A \) in \( \cat{S} \), and if \( \cat{S} \) is skeletal, this object \( A \) is unique.
\end{defproof}

\begin{example}\label{ex:def:category_skeleton}
  We list examples of \hyperref[def:category_skeleton]{skeletons of categories}:
  \begin{thmenum}
    \thmitem{def:category_skeleton/cardinals} The \hyperref[def:small_set]{small} \hyperref[def:cardinal]{cardinal numbers} are a \hyperref[def:category_skeleton]{skeleton} of \( \cat{Set} \).

    Indeed, for every set \( A \), its \hyperref[thm:cardinality_existence]{cardinality} is the unique cardinal isomorphic to \( A \).
  \end{thmenum}
\end{example}

\begin{theorem}[Category skeleton existence]\label{thm:category_skeleton_existence}
  Every \hyperref[def:category]{category} has a \hyperref[def:category_skeleton]{skeleton}.
\end{theorem}
\begin{comments}
  \item In \hyperref[def:zfc]{\logic{ZF}}, this theorem is equivalent to the \hyperref[def:zfc/choice]{axiom of choice}. We list equivalent results in \cref{thm:axiom_of_choice_equivalences}.
\end{comments}
\begin{proof}
  \ImplicationSubProof[def:zfc/choice]{the axiom of choice}[thm:category_skeleton_existence]{skeleton existence} Suppose that the axiom of choice holds. Fix a category \( \cat{C} \).

  Denote by \( \obj(\cat{C}) / \cong \) the quotient of \( \obj(\cat{C}) \) by the isomorphism relation. Using the axiom of choice, we can obtain a \hyperref[def:choice_function]{choice function} \( c: (\obj(\cat{C}) / \cong) \to \obj(\cat{C}) \).

  Define \( \cat{S} \) as the subcategory induced by the set-theoretic image of \( c \). By definition, every object \( A \) in \( \cat{C} \) is isomorphic to \( c(A) \) and to no other object in \( \cat{S} \).

  \ImplicationSubProof[thm:category_skeleton_existence]{skeleton existence}[def:zfc/choice]{the axiom of choice} Suppose that every category has a skeleton.

  Let \( \mscrA \) be a family of nonempty sets. Construct a category \( \cat{C} \) from the \hyperref[def:disjoint_union]{disjoint union} \( \coprod_{A \in \mscrA} A \), where a morphism exists only between members of the same set. This construction is performed in detail in our proof of \fullref{thm:categorical_equivalence_characterization}.

  Then \( \cat{C} \) has a skeleton \( \cat{S} \). All morphisms in \( \cat{C} \) are isomorphisms, hence the set \( \obj(\cat{S}) \) contains exactly one representative for each set in the family \( \mscrA \).

  More precisely, define the set
  \begin{equation*}
    S \coloneqq \set{ x \given \iota_A(x) \in \obj(\cat{S}) }.
  \end{equation*}

  Then \( S \) satisfies \cref{thm:axiom_of_choice_equivalences/choice_sets}.

  Since the family \( \mscrA \) is arbitrary, we conclude that the axiom of choice holds.
\end{proof}

\begin{remark}\label{rem:skeletons_and_preorder_categories}
  Rather than defining representatives of equivalence classes, as in \fullref{thm:category_skeleton_existence}, we can define morphisms between the equivalence classes themselves, as in \cref{def:antisymmetric_quotient}.

  This does not require the \hyperref[def:zfc/choice]{axiom of choice}, but is rarely applicable, unfortunately. One case where it is applicable is in \hyperref[def:thin_category]{thin categories} --- see \fullref{thm:order_category_isomorphism}.
\end{remark}

\paragraph{Monoid delooping}

\begin{definition}\label{def:groupoid}\mcite[20]{MacLane1998CategoryTheory}
  A \term[ru=группоид (\cite[30]{ЦаленкоШульгейфер1974ОсновыТеорииКатегорий})]{groupoid} is a category whose only morphisms are \hyperref[def:morphism_invertibility/isomorphism]{isomorphisms}.
\end{definition}
\begin{comments}
  \item The term \enquote{groupoid} is also sometimes used for sets with binary operations --- see \cref{rem:magma_terminology}. We will avoid the later terminology.
\end{comments}

\begin{definition}\label{def:monoid_delooping}\mcite[def. 1.1.7; 10]{Perrone2024StartingCategoryTheory}
  Let \( (M, \cdot, e) \) be a \hyperref[def:monoid]{monoid}. We define the \term{delooping} \( \cat{B}(M) \) of \( M \) as the following category:
  \begin{thmenum}
    \thmitem{def:monoid_delooping/objects} The only \hyperref[def:category/objects]{object} is \( \Anon \), some (preferably small) set not in \( M \).

    \thmitem{def:monoid_delooping/morphisms} The \hyperref[def:category/morphisms]{set of morphisms} is the underlying set \( M \) of the monoid, with \( \Anon \) as the domain and codomain.

    \thmitem{def:monoid_delooping/composition} The \hyperref[def:category/composition]{composition} \( y \bincirc x \) of \( x \) and \( y \) is their product \( y \cdot x \).

    In accordance with \cref{rem:category_composition_order}, we reverse the order of morphisms when writing their composition. In accordance with \cite[def. 1.1.7]{Perrone2024StartingCategoryTheory}, we use the same order for composition as for multiplication.

    \thmitem{def:monoid_delooping/identity} The \hyperref[def:category/identity]{identity morphism} on \( \Anon \) is \( e \)
  \end{thmenum}
\end{definition}
\begin{comments}
  \item The object \( \Anon \) generally depends on \( M \), but the difference can safely be ignored.
  \item Delooping can act as a functor \( \cat{B}: \cat{Mon} \to \cat{Cat} \) --- see \cref{thm:delooping_as_functor}.
\end{comments}

\begin{proposition}\label{thm:delooping_as_functor}
  \hyperref[def:monoid_delooping]{Delooping} extends to an \hyperref[def:functor_invertibility/injective]{embedding} \( \cat{B}: \cat{Mon} \to \cat{Cat} \) by sending each homomorphism \( \varphi: M \to N \) to the functor whose restriction to morphisms in \( \varphi \).
\end{proposition}
\begin{proof}
  \SubProofOf*{def:functor/CF1} For the identity homomorphism \( \id_M: M \to M \), for each element \( x \) of \( M \), we have
  \begin{equation*}
    \cat{B}(\id_M)(x) = \id_M(x) = \id_{\cat{B}(M)}(x).
  \end{equation*}

  \SubProofOf*{def:functor/CF2} Similarly, for any pair of composable homomorphisms \( \varphi: M \to N \) and \( \psi: N \to K \), for every element \( x \) of \( M \) we have
  \begin{equation*}
    \cat{B}(\psi \bincirc \varphi)(x) = [\psi \bincirc \varphi](x) = [\cat{B}(\psi) \bincirc \cat{B}(\varphi)](x).
  \end{equation*}

  \SubProof{Proof that \( \cat{B} \) is injective on morphisms} Suppose that \( \cat{B}(\varphi) = \cat{B}(\varphi') \) for \( \varphi: M \to N \) and \( \varphi': M' \to N' \).

  Then \( \cat{B}(M) = \cat{B}(M') \) and \( \cat{B}(N) = \cat{B}(N') \). By definition, \( \cat{B}(M) \) has as its morphisms the set \( M \), thus \( M = M' \) and \( N = N' \).

  Moreover, for every element \( x \) of \( M \), we have
  \begin{equation*}
    \varphi(x) = \cat{B}(\varphi)(x) = \cat{B}(\varphi')(x) = \varphi'(x),
  \end{equation*}
  thus \( \varphi \) and \( \varphi' \) coincide as functions.
\end{proof}

\begin{proposition}\label{thm:delooping_of_opposite_monoid}
  The \hyperref[def:monoid_delooping]{delooping} of the \hyperref[def:monoid/opposite]{opposite monoid} of \( M \) is the \hyperref[def:opposite_category]{opposite category} of the delooping of \( M \).
\end{proposition}
\begin{proof}
  Trivial.
\end{proof}

\begin{proposition}\label{thm:delooping_of_group}
  The \hyperref[def:monoid_delooping]{delooping} of a \hyperref[def:group]{group} is a \hyperref[def:groupoid]{groupoid}.
\end{proposition}
\begin{comments}
  \item There is a restricted form of a converse --- see \cref{thm:connected_delooping}.
\end{comments}
\begin{proof}
  Trivial.
\end{proof}

\begin{definition}\label{def:connected_category}\mcite[88]{MacLane1998CategoryTheory}
  We say that a \hyperref[def:category]{category} is \term{connected} its underlying multigraph is \hyperref[def:graph_connectedness/weak]{weakly connected}.
\end{definition}
\begin{comments}
  \item Note that connected \hyperref[def:thin_category]{thin categories} do not correspond to \hyperref[def:binary_relation/connected]{connected relations} (i.e. \hyperref[def:totally_ordered_set]{totally ordered sets}) under \fullref{thm:order_category_isomorphism}. This is simply a terminological clash.
\end{comments}

\begin{proposition}\label{thm:connected_delooping}
  For every \hyperref[def:connected_category]{connected} \hyperref[def:groupoid]{groupoid} \( \cat{G} \) there exists a \hyperref[def:group]{group} \( G \) such that \( \cat{G} \) is \hyperref[def:equivalence_of_categories]{equivalent} to the \hyperref[def:monoid_delooping]{delooping} \( \cat{B}(G) \).
\end{proposition}
\begin{comments}
  \item See \cref{thm:delooping_of_group} for a much simpler converse.
\end{comments}
\begin{proof}
  \SubProof{Proof that all objects are isomorphic} Fix two objects \( A \) and \( B \) in \( \cat{G} \). Since \( \cat{G} \) is connected, there exists an undirected path connecting \( A \) with \( B \). Since \( \cat{G} \) is a groupoid, all morphisms in the path are invertible, and hence there also exists a directed path from \( A \) to \( B \). This directed path can be composed into an isomorphism from \( A \) to \( B \).

  Therefore, all objects in \( \cat{G} \) are isomorphic.

  \SubProof{Construction of group} \Fullref{thm:category_skeleton_existence} implies that there exists a \hyperref[def:skeletal_category]{skeleton} \( \cat{S} \) of \( \cat{G} \).

  Since all objects in \( \cat{G} \) are isomorphic, \( \cat{S} \) can only have a single object, \( \Anon \). Hence, it is possible to compose any two morphisms in \( \cat{S} \).

  Denote by \( G \) be set of morphisms of \( \cat{S} \). Then \( (G, {\bincirc}, \id_{\Anon}) \) is a group.

  \SubProof{Proof of equivalence} Note that the skeleton \( \cat{S} \) and the delooping \( \cat{B}(G) \) are equal. Since \( \cat{G} \) and \( \cat{S} \) are equivalent, so are \( \cat{G} \) and \( \cat{B}(G) \).
\end{proof}
